% Generated draft: requires mathematical and source review.
\documentclass[11pt]{article}
\usepackage[margin=25mm]{geometry}
\usepackage{fontspec}
\setmainfont{Latin Modern Roman}
\usepackage{amsmath,amssymb,mathtools,unicode-math}
\setmathfont{Latin Modern Math}
\usepackage{enumitem,needspace,xurl,hyperref}
\setlength{\parindent}{0pt}
\setlength{\parskip}{5pt}
\setlength{\emergencystretch}{3em}
\newcommand{\N}{\mathbb{N}}
\newcommand{\Z}{\mathbb{Z}}
\newcommand{\Q}{\mathbb{Q}}
\newcommand{\R}{\mathbb{R}}
\newcommand{\C}{\mathbb{C}}
\newcommand{\F}{\mathbb{F}}
\newcommand{\A}{\mathbb{A}}
\newcommand{\PP}{\mathbb{P}}
\newcommand{\E}{\mathbb{E}}
\newcommand{\eps}{\varepsilon}
\newcommand{\dd}{\,\mathrm{d}}
\begin{document}
\section*{UnsolvedMath ID: 20000371}
\textbf{Source problem code:} catalogue entry
\subsection{Definitions and mathematical statement}
The original asks whether Bhargava's arithmetic-statistical methods can yield Hasse-principle results for the following special surfaces and extend Swinnerton-Dyer's diagonal analysis. A Hasse-principle statement relates rational points to points over every completion of $\Q$, with any further hypotheses explicitly required; no universal unconditional local-to-global assertion is stated here.

For homogeneous binary cubics $F,G\in\Q[x,y]$, consider
\[
 S=\{F(x,y)=G(z,w)\}\subset\PP^3_{\Q},\qquad
 u^3=tF(x,y),\quad v^3=tG(z,w).
\]
The latter are common cubic twists, with $t\in\Q^\times$ considered up to cubes when using twist classes. Smoothness/separability hypotheses making these plane curves genus one are not given. The source calls the descriptions equivalent; the report stresses that zero-value boundary cases must be treated separately rather than dividing blindly by $u,v$ or the cubic values.

For separable quartics $f,g$, the second model is
\[
 z^2=f(x)g(y),\qquad u^2=tf(x),\quad v^2=tg(y),
\]
with common quadratic twist class $t\in\Q^\times/\Q^{\times2}$. The source calls the surface Kummer K3. A proper smooth model, affine boundary cases, and the distinction between affine quartic points and points on projective genus-one torsors must be specified before making rational-point equivalences. Projective points at infinity alone do not automatically supply the required affine solution.

The stated goal is to remove the finiteness assumption on the $2$-primary Tate--Shafarevich group of the relevant quadratic twists. Here $\operatorname{Sha}(E/\Q)$ is the kernel of $H^1(\Q,E)\to\prod_vH^1(\Q_v,E)$ for the relevant Jacobian elliptic curve; the extracted symbol ``X'' is interpreted in that sense by the research report. Which twists and auxiliary descent hypotheses occur in Swinnerton-Dyer's argument are not supplied.

The report attributes positive-proportion rational-point results for split cubic surfaces, and analogous Kummer families, to Browning, while explicitly leaving general fixed-surface and finiteness-removal questions open. It proposes exact common-twist/value-class criteria, including cubic zero cases and the affine/projective distinction, but does not include the complete criteria.

Its relative selector requires exceptional estimates within the \emph{same} admissible common twist family. For a specified height ordering and admissible sets $\mathcal A(T)$, a sufficient proposed condition is that the sum of the two marked-class bad counts is $o(|\mathcal A(T)|)$, so some admissible twist is good for both coverings. Density-zero exceptions in the ambient unrestricted family do not imply this relative estimate for a sparse $\mathcal A$. The admissible conditions, marked classes, and ordering are not defined. These are unreviewed reductions and family-average reports, not a solution for every fixed surface.
\subsection{Short English statement}
Use arithmetic statistics to find one common twist with the required rational points and improve local-to-global arguments for the special cubic and Kummer surfaces. The report distinguishes family-average progress from fixed-surface results and insists on relative exceptional estimates inside the common admissible twist family. The original finiteness-removal objective remains open in its assessment.
\subsection{Sources}
\begin{itemize}
\item Local UnsolvedMath registry, permanent ID $20000371$, source code \texttt{catalogue entry}; created and updated 2026-08-14. Status \texttt{partially\_solved}; classification \texttt{PARTIAL-PROGRESS}; difficulty L3; proposer field null, while source attribution names Alexei Skorobogatov.
\item AIM open-problem notes, \emph{Rational and integral points on higher-dimensional varieties}, Problem 3, dated May 28, 2014, source type PDF, \url{https://aimath.org/pastworkshops/ratlhigherdimvarproblems.pdf}. Canonical location \texttt{aim-algebraic-number-theory-notes.json}, \texttt{notes[22]}, tag \texttt{problem}; statement status \texttt{exact}. Extracted page/header text is not part of the mathematical conditions. Original remarks and literature field empty.
\item Original reference: Swinnerton-Dyer, ``[Ann. Sci. ENS],'' on the diagonal search for common twists, without title/year. The source's ``recent'' Bhargava talk belongs to its 2014 setting. The research summary names Browning for positive-proportion results without exact works or height/family conventions.
\item Supplied research attempt 1: result \texttt{partial\_result}, run status \texttt{partially\_solved}, rigor label \texttt{proved}, low novelty confidence. Candidate contribution \texttt{reduction}: common-twist/value criteria and a relative marked-torsor selector, sharp against ambient-density substitution. Explicit caveat: machine-generated research aid, not peer reviewed; independently verify before citation.
\item Dataset attribution: Ulam AI / UnsolvedMath Contributors, CC BY 4.0; original problem attribution remains to the supplied AIM workshop notes.
\end{itemize}
\end{document}
